\section{Conclusion}
The Alkhidmat Public Chat Portal successfully bridges the communication gap between one of Pakistan's largest NGOs and its stakeholders. By transitioning from manual query handling (phone/email) to an AI-driven, multilingual RAG system, the project has achieved several key milestones:
\begin{itemize}
    \item \textbf{Operational Efficiency:} Repetitive queries regarding blood bank locations and donation workflows are now handled autonomously, reducing staff burden.
    \item \textbf{Inclusivity:} The support for Roman Urdu and Urdu script makes the platform accessible to low-literacy users and the overseas diaspora.
    \item \textbf{Trust and Transparency:} The integration of Self-RAG verification and explicit source attribution ensures that the information provided to donors and beneficiaries is verified and grounded in official SOPs.
\end{itemize}

\section{Future Work}
While the current system provides a robust MVP (Phase 1), the following enhancements are proposed for future development:
\begin{itemize}
    \item \textbf{Optimized Response Time:} Improving the system's response time through optimization of the RAG pipeline and caching strategies, so system can answer queries in under 10 seconds for most cases.
    \item \textbf{Bilingual Voice Integration:} Implementing Speech-to-Text (STT) capabilities to allow rural beneficiaries to interact with the bot via voice notes in Urdu and regional dialects.
    \item \textbf{Real-time Inventory Access:} Moving beyond periodic updates to direct API integration with Alkhidmat’s healthcare and blood bank management systems for live stock availability.
    \item \textbf{Expanded Language Support:} Extending the NLP pipeline to include regional languages such as Sindhi, Pashto, and Punjabi to enhance localized outreach.
    \item \textbf{Advanced Analytics:} Implementing predictive modeling on the Admin Dashboard to forecast donation trends and healthcare service demand based on user query patterns.
\end{itemize}